\section{自动驾驶软件栈：从规则到云端模型工厂}

上一章讲人和组织，本章进入技术栈。刘先明把自动驾驶演进拆成几个阶段：Software 1.0 是规则和优化；Software 1.5 是模型加规则；Software 2.0 是端到端；Software 3.0 或 VLA/VLM 路线则是在更大算力和更大数据下，用更大模型做数据 scaling。最终，车端硬件无法承载全部训练和大模型，于是需要云端模型工厂。

\begin{figure}[H]
\centering
\includegraphics[width=0.96\textwidth]{autonomous-stack-evolution.png}
\caption{自动驾驶软件栈演化：从规则、半模型、端到端，到更大模型和云端工厂。自制概念图，依据 00:02:16--00:19:00 对谈内容整理。}
\end{figure}

\begin{knowledgebox}{读图：每一代都在拆掉旧上限}
Software 1.0 被规则和优化约束；Software 1.5 的上限仍在规控代码；Software 2.0 用端到端减少人工结构；Software 3.0/VLA/VLM 用更大模型和更多数据继续 scaling；云端工厂负责训练、蒸馏、量化、剪枝和部署。
\end{knowledgebox}

\subsection{Software 1.0 到 2.0}

Software 1.0 使用激光雷达聚类、传统检测、手写规则和数学优化来做感知与规控。这一代系统可解释、工程可控，但会被规则上限卡住。Software 1.5 把模型引入感知，例如检测和分割，但规划控制仍依赖大量规则。Software 2.0 则试图让神经网络从数据中学习更完整的输入输出关系，用数据迭代代替越来越复杂的手写规则。

\begin{warningbox}{规则不是错，规则是阶段性上限}
规则系统在早期非常重要，因为它可控、可调、可解释。但当场景复杂度超过人工规则维护能力时，规则会变成上限。拆规则不是否定工程，而是承认规模化问题需要数据驱动。
\end{warningbox}

\subsection{云端模型工厂}

车端芯片无法直接训练超大模型，也不适合直接运行所有云端能力。因此，刘先明提出“云端工厂”思路：在云端训练巨大模型，再通过蒸馏、量化、剪枝等方式压缩能力，部署到不同车端硬件上。这个模式类似模型生产线，训练好一个云端母模型后，可以持续生成适配不同平台的车端模型。

\begin{figure}[H]
\centering
\includegraphics[width=0.96\textwidth]{cloud-model-factory.png}
\caption{云端模型工厂：大模型在云端训练，再蒸馏、量化、剪枝到车端。自制概念图，依据 00:02:16--00:19:00 对谈内容整理。}
\end{figure}

\begin{knowledgebox}{读图：云端负责能力，车端负责部署}
云端大模型吃大参数和大数据，训练后通过蒸馏把能力迁到小模型，再量化剪枝适配车端算力。这样既能利用 scaling，又能满足量产车硬件约束。
\end{knowledgebox}

\subsection{本章小结}

自动驾驶软件栈演化的核心，是不断拆掉人工中间层，扩大数据和模型规模，再用工程体系把云端能力压缩到车端。云端模型工厂是连接大模型和量产部署的关键结构。

